\chapter{Glossary and Notation Index}
\label{appD:chapter}

This appendix collects every technical term and symbol used in the
monograph, with a short self-contained definition and a plain-text
pointer to the chapter or appendix where the notion is developed. It is
a reading aid, not a normative source: where this glossary and a
definition in the body could be read differently, the body governs, and
where the body and a registered ledger artifact could be read
differently, the ledger governs. Chapter references are given as plain
text (for example ``Chapter 10'') rather than as links.

\section{Terms}
\label{appD:sec:terms}

\textbf{ACBE mark.} The first coordinate of an atom of $\Omega$, a
formal label preserved for controlled migration vocabulary from the
legacy construction. It carries no empirical interpretation until a
separate operationalization and validation record exists; no theorem in
the monograph attaches meaning to it beyond its combinatorics. Standing
data in Chapter 10; provenance role in Chapter 11.

\textbf{Admission gates.} The ten conjunctive conditions --- independent
review, formalization, typed product-site covers, the product-site
lift, the backward-compatibility proofs, the cross-context decision,
probabilistic gates, novelty search, registry and IP review, external
referee review --- that must all close before Noetic Topos 2.0 may
leave Ring 4. Enumerated and discussed in Chapter 12.

\textbf{Admission ring (Ring 4).} The program's containment level for
the entire artifact set of this monograph: internal research material,
not licensed for public, empirical, or product claims. Stated in the
status boundaries of Chapters 10, 11, and 12.

\textbf{Ambient topos.} The presheaf topos
$\Eamb = \PSh(\cC \times L)$ over the product of the legacy category
with the context lattice; the carrier of the three modalities and the
home of unresolved and partially resolved evidence. The candidate
definition sets $\mathrm{NoeTop}_2 := \Eamb$. Chapter 12.

\textbf{Associated sheaf (sheafification).} The universal sheaf
receiving a map from a given presheaf: the left adjoint to the
inclusion of sheaves into presheaves. On the finite context factor it
is realized exactly by the one-partition reflector. Chapter 10; general
background in the cited Stacks Project material.

\textbf{Atom.} One of the four elements
$\omega_{ij} = (i, j)$ of $\Omega$, carrying a formal ACBE mark $i$
and RPM mark $j$. Chapter 10; tabulated in Appendix A.

\textbf{Backward-compatibility conjecture.} The unproved comparison
$\Eres \simeq \Eleg^4$ together with the neutral-channel embedding
$\iota$ and the natural isomorphism
$\mathrm{ev}_{00} \circ \iota \cong \mathrm{Id}$; ledger status
\textsc{proof required}. Two round trips (categorical and strict
serialization) are required and neither is complete. Chapter 12.

\textbf{Bit-assignment presheaf.} The value shape of the witness
presheaves: functions from a set of complete blocks to $\{0, 1\}$,
restriction forgetting coordinates. Also the name of an executable
fixture in the composition verifier. Appendix C; Chapter 11.

\textbf{Block.} A member of one of the partitions of $\Omega$: an
ACBE block $A_i$, an RPM block $R_j$, or a singleton block of $Q$.
Chapter 10; Appendix A.

\textbf{Certificate (finite model).} The registered artifact
\texttt{FINITE\_MODEL\_CERTIFICATE\_v0.1.json}: the expected-result
record of the finite-site checker, comprising the twenty-three
registered checks and the aggregate pass flag, under its own claim
boundary. Chapter 11.

\textbf{Claim boundary.} A first-class statement, attached to an
artifact, of what the artifact does not establish. The monograph's
governing example is the certificate's boundary: exhaustive execution
is not proof-assistant verification, peer review, empirical
validation, or a product-claim license. Chapters 10, 11, 12 and the
appendices each carry their own.

\textbf{Coincidence (of unmarked families).} The phenomenon that
$K_A(V)$ and $K_R(V)$ can be equal as bare arrow sets --- at the empty
context, the four singletons, and the two diagonals --- so that the
topology retains no trace of which partition produced the cover, while
the marked records remain distinct. Derived in Appendix A; singleton
cases registered in the certificate, Chapter 11.

\textbf{Common refinement.} The family of pairwise intersections of
two covers. For the two global covers it is the four-singleton cover,
by transversality; as a marked record it is a derivation node carrying
both parent digests. Chapters 10 and 11.

\textbf{Complete block.} A block $B$ with $B \subseteq V$ for the
context $V$ at hand; the witness presheaves assign bits only to
complete blocks. Appendix A tabulates completeness at every context.

\textbf{Composition theorem.} The chapter-10 result
$M_A M_R \cong M_{\mathrm{atom}} \cong M_R M_A$, unit-compatibly and
naturally, with either composite naturally isomorphic to
sheafification for $J_A \vee J_R$. Status: internal proof candidate,
review required. Chapter 10; executable corroboration in Chapter 11.

\textbf{Context (context lattice).} An element of
$L = \mathcal{P}(\Omega)$, ordered by inclusion; the sixteen-object
poset carrying the finite theory. Chapter 10; enumerated in Appendix
A.

\textbf{Covering sieve.} A sieve belonging to the topology at its
object; the sheaf condition quantifies over covering sieves. Chapter
10; counted exactly in Chapter 11 and Appendix A.

\textbf{Cover kind.} The field of a mark record naming the admitting
clause (\textsc{observational}, \textsc{acbe}, or \textsc{rpm});
distinct cover kinds at coinciding unmarked families are what the
marked layer preserves. Chapter 11.

\textbf{Dedekind number.} The number $M(n)$ of down-sets of the
Boolean lattice on $n$ points ($2, 3, 6, 20, 168$ for
$n = 0, \ldots, 4$); used to derive the registered total of $298$
sieves. Chapter 11; per-context table in Appendix A.

\textbf{Derivation kind.} The field of a marked cover declaring how it
arose: \textsc{generator}, \textsc{pullback}, or
\textsc{common refinement}; derived kinds require parent resolution.
Chapter 11.

\textbf{Digest.} The SHA-256 hash of a marked cover's canonical JSON,
binding target, members, mark, and sorted parent digests; the unit of
provenance chaining. Chapter 11.

\textbf{Distinctness theorem.} Theorem NT2-3: $F_A$ is a $J_A$-sheaf
but not a $J_R$-sheaf and symmetrically, so the two sheaf
subcategories differ and neither modality is trivial; scoped to
$\PSh(L, \Set)$. Witness tables in Appendix C; verdict checks in
Chapter 11.

\textbf{Empty-product convention.} The clause
$(R_P F)(\varnothing) = \mathbf{1}$: the reflector's value at the
empty context is the singleton regardless of $F(\varnothing)$, forced
by the null cover. Load-bearing in the sheaf lemma and exercised by a
dedicated adversarial fixture. Chapter 10; Chapter 11.

\textbf{Evidence Packet.} The enclosing program artifact that must
separately bind raw ingress bytes; the cover digest is the canonical
semantic-record digest only, not a substitute for that binding.
Chapter 11 context; program-level notion.

\textbf{Expressivity boundary.} The observation that, even if the
backward-compatibility conjecture is proved, $\Eleg^4$ is four
independent channels with no cross-atom coupling, so the finite model
is a separation and provenance reference model rather than a final
semantics. Chapter 12.

\textbf{Fail-closed.} The integrity posture of the marked layer:
unknown or duplicate JSON keys, uncontrolled review states, and
derived records with wrong or unresolvable parents are rejected
rather than treated as evidence. Chapter 11.

\textbf{Flattening and regrouping.} The mutually inverse bijections
$\varphi$ and $\psi$ between the nested product
$(R_A R_R F)(V)$ and the atomic product $(R_Q F)(V)$, induced by
transversality; the computational core of the composition theorem.
Chapter 10.

\textbf{Gluing (matching family).} A matching family for a cover
assigns a compatible local section to each member; the sheaf
condition asks that the canonical map from global sections to
matching families be a bijection. Chapter 10; computed for the
witnesses in Appendix C.

\textbf{Grothendieck topology.} An assignment of covering sieves to
objects satisfying maximality (J1), pullback stability (J2), and
local character (J3). Chapter 10; computed exactly by saturation in
Chapter 11.

\textbf{Join (of topologies).} The least topology containing two
given topologies. The monograph's central join identification is
$J_A \vee J_R = J_{\mathrm{atom}}$. Chapter 10.

\textbf{Lean layer.} The bounded proof-assistant scope: four-atom
exhaustiveness, proper global pieces, the transverse refinement
identity, and the witness-map lemmas --- and nothing else; the sheaf
theory is explicitly outside it. Chapter 11.

\textbf{Legacy collapse.} The diagnosis that the legacy ACBE and RPM
cover clauses generate only maximal sieves (identity via reflexivity;
identity via split-epi closure), so the legacy topology reduces to
its jointly-surjective clause and retains no modal provenance.
Executable checks in Chapter 11; motivating result of the baseline
chapters.

\textbf{Legacy site.} The registered category
$\cC = \mathrm{Noetica}$ with its jointly-surjective topology
$J_{\mathrm{legacy}}$ (also written $J_{\mathrm{obs}}$ in the ambient
package) and sheaf topos $\Eleg$. Chapter 12.

\textbf{LP feasibility certificate.} The second probabilistic gate:
existence of a global joint distribution over a fixed finite cover
posed as a finite linear program, with a feasibility or infeasibility
certificate as the licensed artifact. Chapter 12.

\textbf{Marginal problem.} The classical fact that
pairwise-compatible marginal distributions need not extend to a joint
distribution; the reason Set-valued gluing licenses no probabilistic
claim. Chapter 12, with the cited classical and sheaf-theoretic
literature.

\textbf{Marked cover.} The immutable record (derivation kind, target,
exact member arrows, mark record, parent digests) that travels in
parallel with the unmarked family; the topology consumes the family,
governance consumes the record. Chapter 11.

\textbf{Mark record.} The provenance component of a marked cover:
cover kind, source authority reference, assumptions, evidence
references, review state, provenance chain. Chapter 11.

\textbf{Maximal sieve.} The sieve of all arrows into an object;
covering in every topology. A sieve containing an identity, or a
split epimorphism with its section, is maximal --- the two collapse
mechanisms. Chapters 10 and 11.

\textbf{Modality.} A left-exact idempotent monad on a presheaf topos;
here, sheafification-followed-by-inclusion for a topology refinement
($M_A$, $M_R$, $M_{\mathrm{atom}}$, $M_{\mathrm{obs}}$). Trivial
means naturally isomorphic to the identity. Chapter 10; ambient roles
in Chapter 12.

\textbf{Modality-composition verifier.} The second executable layer:
implements the explicit reflector, its unit, and both comparison
maps, and computes both topologies exactly by sieve enumeration and
saturation. Chapter 11.

\textbf{Model ID.} The registered identifier
\texttt{NT2-CONTEXT-2X2-TRANSVERSE-PARTITIONS} naming the finite
model checked by the certificate. Chapter 11.

\textbf{Neutral channel.} The registry-fixed atom $\omega_{00}$ used
by the conjectural legacy embedding $\iota : X \mapsto (X, 1, 1, 1)$;
evaluation there is $\mathrm{ev}_{00}$. Chapter 12.

\textbf{NoeSystem ($\mathrm{NoeSystem}_2$).} The full research
object --- the tuple of
$\mathrm{NoeTop}_2$, the three modalities, the resolved subtopos, and
the external mark record. A structured package, not a topos or
category. Chapter 12.

\textbf{NoeTop ($\mathrm{NoeTop}_2$).} The underlying topos
candidate, defined as the
ambient presheaf topos $\Eamb$ (not the resolved subtopos, on which
the modalities restrict to identity). Chapter 12.

\textbf{Novelty status.} All theorem-ledger records carry
\textsc{novelty unassessed}: no dedicated literature search has been
performed, and the monograph does not call its results new theorems
in the scholarly sense. Chapter 12.

\textbf{Null cover.} The empty family at the empty context, which
covers it in all three bases; forces every sheaf to have a
one-element value at $\varnothing$. Chapter 10.

\textbf{Obligation (cross-context semantics).} The registered record
\textsc{obl-nt2-cross-context-semantics}: choose and formalize either
the ambient modal package as principal semantics or a resolved
semantics with typed overlap or transition data and proved coherence,
under four acceptance criteria. Chapter 12.

\textbf{One-partition reflector.} The explicit construction
$(R_P F)(V) = \prod_{B \in I_P(V)} F(V \cap B)$ with the empty
product a singleton, restriction dropping absent coordinates, and
unit given by restrictions; proved to be the associated-sheaf
reflector for $J_P$. Chapter 10.

\textbf{Presheaf.} A functor $L^{\mathrm{op}} \to \Set$ (more
generally $\cC^{\mathrm{op}} \to \Set$): sets of sections with
compatible restriction maps. Chapter 10.

\textbf{Proof-obligation ledger.} The registered artifact recording
open design-and-proof obligations with statements, dependencies,
acceptance criteria, and claim boundaries. Chapter 12.

\textbf{Provenance chain.} The parent-digest lineage of a marked
cover, checked acyclically and recomputed on resolution. Chapter 11.

\textbf{Pullback sieve; pullback stability.} The pullback of a sieve
along an arrow ($g^{*}S$); axiom (J2) requires covering sieves to
pull back to covering sieves, and the bases satisfy it at the family
level (in $L$, pullback is intersection). Chapters 10 and 11.

\textbf{Reflector.} A left adjoint to the inclusion of a reflective
subcategory; here always the associated-sheaf functor for one of the
topologies, realized explicitly by $R_P$. Chapter 10.

\textbf{Representable presheaf.} $yW = \Hom_L(-, W)$: a singleton at
$V \subseteq W$, empty otherwise. All sixteen are sheaves for both
factor topologies (subcanonicity). Chapter 10 usage; check in
Chapter 11.

\textbf{Resolved subtopos.} $\Eres = \Sh(\cC \times L,
J_{\mathrm{res}})$, the sheaf subcategory for the fully joined
refinement; modally inert by the resolved-identity theorem. Chapters
10 and 12.

\textbf{Restriction map.} The action of a presheaf on an inclusion
$U \subseteq V$, written $x \mapsto x|_U$; for the witnesses,
coordinate forgetting (Appendix C).

\textbf{RPM mark.} The second coordinate of an atom; formal label
only, as for the ACBE mark. Chapter 10.

\textbf{Section-retraction.} A pair
$\eta : X \to R(X)$, $\varepsilon : R(X) \to X$ with
$\varepsilon \circ \eta = \id$; the legacy RPM data. Its split
epimorphism is one of the two collapse mechanisms. Chapter 11.

\textbf{Separatedness.} The injectivity half of the sheaf condition:
sections agreeing on a cover are equal. Used repeatedly in the
universal-property proof. Chapter 10.

\textbf{Sheaf.} A presheaf whose gluing map is a bijection for every
covering sieve; equivalently, for a coverage basis, for every basis
family. Chapter 10.

\textbf{Sieve.} A precomposition-closed set of arrows into a fixed
object; in the poset $L$, a down-set of subcontexts. Chapter 10;
counted in Chapter 11 and Appendix A.

\textbf{Singleton-factor bijection.} The canonical identification
$(R_P F)(W) = F(W)$ at a nonempty context $W$ contained in a single
block; the engine of the uniqueness half of the universal property.
Chapter 10.

\textbf{Site.} A category with a Grothendieck topology. Chapter 10.

\textbf{Split epimorphism.} An arrow with a section; a sieve
containing one is maximal. Chapters 10 and 11.

\textbf{Subcanonical.} A topology for which every representable
presheaf is a sheaf; both factor topologies are subcanonical, while
legacy subcanonicity is not derivable. Chapter 10 usage; Chapter 12
consequences.

\textbf{Theorem ledger.} The registered record set assigning each
theorem its proof status; drafting prose promotes nothing recorded
there. Status recaps in Chapters 10 and 12.

\textbf{Transverse partitions.} The two partitions
$\{A_0, A_1\}$ and $\{R_0, R_1\}$ of $\Omega$ with
$A_i \cap R_j = \{(i, j)\}$; transversality is the single identity
powering the join and composition results. Chapter 10.

\textbf{Unit.} The natural transformation $\eta_P : \id \Rightarrow
R_P$ (restrictions to pieces), universal among maps to sheaves;
unit compatibility is part of the composition theorem. Chapter 10.

\textbf{Universal property.} The bijection
$\Hom_{\Sh}(R_P F, G) \cong \Hom_{\PSh}(F, G)$ natural in both
variables, exhibiting $R_P$ as THE associated-sheaf reflector.
Chapter 10.

\textbf{Witness presheaf.} $F_A$ or $F_R$: the explicit bit-assignment
presheaves separating the two sheaf subcategories. Appendix C
tabulates them in full.

\section{Symbols}
\label{appD:sec:symbols}

\begin{longtable}{@{}lll@{}}
\toprule
Symbol & Meaning & Where \\
\midrule
\endhead
\bottomrule
\endfoot
$\Omega$ & the four-atom set $\{0,1\} \times \{0,1\}$ & Chapter 10 \\
$\omega_{ij}$ & the atom $(i, j)$; $\omega_{00}$ is the neutral atom &
Chapter 10, 12 \\
$L$ & the context lattice $\mathcal{P}(\Omega)$ & Chapter 10 \\
$A_i$, $R_j$ & ACBE and RPM blocks & Chapter 10, Appendix A \\
$Q$ & the singleton partition of $\Omega$ & Chapter 10 \\
$\Delta$, $\nabla$ & the diagonal contexts
$\{\omega_{00}, \omega_{11}\}$, $\{\omega_{01}, \omega_{10}\}$ &
Appendix A \\
$T_{ij}$ & the triple $\Omega \setminus \{\omega_{ij}\}$ &
Appendix A \\
$K_P(V)$ & generating family of partition $P$ at $V$ & Chapter 10 \\
$J_A$, $J_R$ & the two factor topologies & Chapter 10 \\
$J_{\mathrm{atom}}$ & the singleton-generated topology; preferred
name $J_{\mathrm{singleton}}$ & Chapter 10 \\
$J_{\mathrm{obs}}$, $J_{\mathrm{legacy}}$ & the legacy
jointly-surjective topology & Chapter 12 \\
$J_{\mathrm{res}}$ & the fully joined refinement
$J_{\mathrm{obs}} \vee J_A \vee J_R$ & Chapter 12 \\
$\vee$ & join of topologies (least upper bound) & Chapter 10 \\
$\le$ & comparison of topologies (containment of covering sieves) &
Chapter 10 \\
$I_P(V)$ & blocks of $P$ meeting $V$ & Chapter 10 \\
$R_P$ & the one-partition reflector for $P$ & Chapter 10 \\
$\eta_P$ & the unit of $R_P$ & Chapter 10 \\
$\varphi$, $\psi$ & flattening and regrouping bijections &
Chapter 10 \\
$M_A$, $M_R$, $M_{\mathrm{atom}}$, $M_{\mathrm{obs}}$ &
sheafification-then-inclusion modalities & Chapters 10, 12 \\
$F_A$, $F_R$ & the witness presheaves & Appendix C \\
$yW$ & the representable presheaf of $W$ & Chapters 10, 11 \\
$\mathbf{1}$ & a one-element set; the empty product & Chapter 10 \\
$\PSh$, $\Sh$ & presheaf and sheaf categories & throughout \\
$\Set$ & the category of small sets & throughout \\
$\cC$ & the registered legacy category $\mathrm{Noetica}$ &
Chapter 12 \\
$\Eamb$, $\Eres$, $\Eleg$ & ambient, resolved, and legacy topoi &
Chapter 12 \\
$\mathrm{NoeTop}_2$, $\mathrm{NoeSystem}_2$ & topos candidate and
full package & Chapter 12 \\
$\iota$, $\mathrm{ev}_{00}$ & conjectural neutral-channel embedding
and evaluation & Chapter 12 \\
$M(n)$ & the Dedekind number (down-sets of the $n$-cube) &
Chapter 11, Appendix A \\
$m_C$ & the maximal sieve on $C$ & Chapter 10 \\
$g^{*}S$ & the pullback sieve of $S$ along $g$ & Chapter 10 \\
$\id$, $\Hom$, $\Ob$ & identity, hom-sets, objects & throughout \\
\end{longtable}

\begin{remark}
\label{appD:rem:statusvocab}
Status vocabulary, for completeness: \textsc{ring 4} (containment
level of the artifact set), \textsc{not licensed} (no public,
empirical, or product claim), \textsc{internal proof candidate /
review required} (an informal proof present, no independent review),
\textsc{proof required} (a conjecture with no proof),
\textsc{novelty unassessed} (no literature search performed), and
\textsc{open} (a registered obligation not yet discharged). These
words are load-bearing wherever they occur; this monograph uses them
in exactly the senses of the registered ledgers, and Chapter 12
collects the gates under which any of them may change.
\end{remark}
